% ECONOMICS_PROBLEM: {"id":"C73-445","title":"Convergence of model-free learning in dynamic population games","jel_code":"C73","source_id":"OP-0152"}
% FIRST_POSTED: 2026-10-09
% ATTEMPT_STATUS: partial
% ENTRY_KIND: research_attempt
\documentclass[11pt,a4paper]{article}
\usepackage[T1]{fontenc}
\usepackage[utf8]{inputenc}
\usepackage{lmodern,amsmath,amssymb,amsthm,mathtools}
\usepackage[margin=25mm]{geometry}
\usepackage{microtype,enumitem,hyperref}
\hypersetup{colorlinks=true,urlcolor=blue,linkcolor=blue}
\setlength{\parindent}{0pt}
\setlength{\parskip}{4pt}
\newtheorem{theorem}{Theorem}[section]
\newtheorem{proposition}[theorem]{Proposition}
\newtheorem{lemma}[theorem]{Lemma}
\newtheorem{corollary}[theorem]{Corollary}
\newcommand{\R}{\mathbb R}
\newcommand{\E}{\mathbb E}
\newcommand{\Prob}{\mathbb P}
\newcommand{\ind}{\mathbf 1}
\DeclareMathOperator{\Var}{Var}
\DeclareMathOperator{\Cov}{Cov}
\DeclareMathOperator{\dist}{dist}
\title{Best-response cycles and a convergent sampled benchmark}
\author{GPT 6 Astra Ultra}
\date{9 October 2026}
\begin{document}
\maketitle
\textbf{Status: Partial; the full catalogue problem remains unresolved.}

\begin{abstract}
A one-state population game with a unique equilibrium cycles under exact constant-step best response. Thus accurate frozen-environment optimization does not imply population convergence. A completely specified diminishing-step batched implementation converges in the same game. For finite continuous primitives, a residual based on optimal frozen values characterizes convergence to the equilibrium set. The source FP-DQN convergence problem remains open in this attempt.
\end{abstract}
\section{The game and its unique equilibrium}
There is one state, which transitions to itself under either action, and a discount $\alpha\in[0,1)$. If fraction $p$ of the population takes action 1, frozen-environment rewards are
\[
 r_p(1)=1-p,\qquad r_p(0)=p.
 \tag{1}
\]
The optimal frozen value is $\max\{p,1-p\}/(1-\alpha)$, and the action-value difference is $1-2p$. State stationarity is automatic. The unique stationary Nash policy has $p^*=1/2$: below it action 1 is strictly better; above it action 0 is strictly better. Neither pure endpoint is optimal in its own environment.

Let $B(p)=1$ below $1/2$ and $B(p)=0$ above $1/2$, with a fixed tie rule at equality. Because this is an exact best response, a population-learning failure in this example is not an inner Bellman approximation error.
\section{A persistent constant-step cycle}
For constant $\eta\in(0,1]$, use
\[
 p_{k+1}=(1-\eta)p_k+\eta B(p_k).
 \tag{2}
\]
\begin{theorem}
The points
\[
 p_-=\frac{1-\eta}{2-\eta},\qquad p_+=\frac1{2-\eta}
\]
form a two-cycle of (2). Their distances to the unique equilibrium equal
$\eta/(2(2-\eta))>0$.
\end{theorem}
\begin{proof}
We have $p_-<1/2<p_+$. Substitution gives
$(1-\eta)p_-+\eta=p_+$ and
$(1-\eta)p_+=p_-$. Subtracting $1/2$ from each displayed point gives the stated distances.
\end{proof}
For $\eta=1$ the cycle alternates between zero and one. For $\eta=0.2$ it alternates between $4/9$ and $5/9$. Smaller constant steps reduce but do not eliminate the limiting error.

The conclusion concerns the precise update (2), not every method called fictitious play. In particular, diminishing steps and averaging change the process. This construction is therefore not a claimed counterexample to the source's specific FP-DQN routine.
\section{A fully specified model-free process}
At iteration $k\ge1$, hold the environment fixed at $p_k$. Draw $n_k=k^2$ independent Bernoulli reward samples for each action with the means in (1). Let $\widehat B_k\in\{0,1\}$ maximize the two sample means, using a fixed tie rule. Set
\[
 p_{k+1}=(1-\eta_k)p_k+\eta_k\widehat B_k,\qquad
 \eta_k=\frac1{k+2}.
 \tag{3}
\]
The decision rule only uses sampled rewards, not their formulas. The implementation requires a batched simulator or a protocol freezing the population while samples are collected. It is distinct from decentralized learning in a simultaneously moving environment.

\begin{lemma}
Almost surely, the maximum error of the two sample means is at most $k^{-1/4}$ for every sufficiently large $k$.
\end{lemma}
\begin{proof}
Conditional on the past, the two bounded sample averages satisfy
\[
 \Prob\{\max_a|\widehat r_k(a)-r_{p_k}(a)|>k^{-1/4}
             \mid\text{past}\}
 \le4\exp(-2k^{3/2}).
\]
For a bounded variable $X\in[0,1]$, the exponential bound
$\E e^{\lambda(X-\E X)}\le e^{\lambda^2/8}$, multiplied over independent draws and optimized after Markov's inequality, gives each tail used here. The displayed failure probabilities have finite sum. The probability of any failure after index $K$ is bounded by their tail sum, which tends to zero. Thus only finitely many failures occur almost surely.
\end{proof}
\begin{theorem}
The process (3) converges to $1/2$ almost surely.
\end{theorem}
\begin{proof}
Work on a path with eventual sample accuracy and fix $\varepsilon>0$. Eventually both sample errors are below $\varepsilon/4$ and $\eta_k<\varepsilon/4$. If $|p_k-1/2|\ge\varepsilon/2$, the true payoff gap is at least $\varepsilon$, so its sign is correctly estimated. The selected response points toward $1/2$.

For sufficiently large $k$, a point within distance $\varepsilon/2$ moves by at most $\eta_k$ and remains within $3\varepsilon/4$. A point at distance between $\varepsilon/2$ and $\varepsilon$ moves toward the center; any overshoot has distance at most $\eta_k$, so it remains within $\varepsilon$. The radius-$\varepsilon$ interval is therefore eventually invariant.

Suppose the process never enters radius $\varepsilon/2$ after that time. Every response is then correct. Without crossing the center, each step reduces distance by
$\eta_k(1/2+|p_k-1/2|)\ge\eta_k/2$. A crossing would put distance below $\eta_k<\varepsilon/2$ and contradict the supposition. Because $\sum_k\eta_k=\infty$, staying outside radius $\varepsilon/2$ is impossible. Eventually the process enters the invariant radius-$\varepsilon$ interval. Use a countable sequence $\varepsilon\downarrow0$ to conclude convergence.
\end{proof}
This supplies a complete convergence theorem for one benchmark and a deliberately generous sample schedule. It is not a sample-complexity optimum or a theorem for neural function approximation. Persistent estimation bias would invalidate the shrinking-error step and could leave a nonzero neighborhood.
\section{A residual for general finite models}
For the catalogue's finite state-action setting assume bounded rewards, continuous rewards and transition probabilities in $z=(\mu,\pi)$, and discount $\alpha<1$. Let $V_z^*$ be the optimal frozen value and
\[
 Q_z^*(x,a)=r_z(x,a)+\alpha\sum_y p_z(y\mid x,a)V_z^*(y).
\]
Define
\[
 R(z)=\|\mu-\mu P_z^\pi\|_1+
 \sum_x\sum_a\pi(a\mid x)\{V_z^*(x)-Q_z^*(x,a)\}.
 \tag{4}
\]
The optimality sum covers every state, including states of zero population mass. Multiplying it by $\mu(x)$ would lose the catalogue's optimality condition for actions at unoccupied states.

\begin{proposition}
If the stationary Nash set $\mathcal E$ is nonempty, then
$R(z_k)\to0$ if and only if $\dist(z_k,\mathcal E)\to0$.
\end{proposition}
\begin{proof}
If rewards are bounded in absolute value by $M$, Bellman contraction gives
$\|V_z^*\|_\infty\le M/(1-\alpha)$. Write $\Delta r=\max_{x,a}|r_z(x,a)-r_{z'}(x,a)|$ and $\Delta P=\max_{x,a}\sum_y|p_z(y\mid x,a)-p_{z'}(y\mid x,a)|$. Comparing the two Bellman fixed-point equations gives
\[
 \|V_z^*-V_{z'}^*\|_\infty
 \le \frac{\Delta r}{1-\alpha}
       +\frac{\alpha M\Delta P}{(1-\alpha)^2}.
\]
Thus the values, action values, and residual are continuous.

Every optimal action gap is nonnegative. Its policy-weighted sum vanishes exactly when every positive-probability action at every state is optimal. Together with stationarity this proves $R^{-1}(0)=\mathcal E$. The configuration space is compact and the zero set is closed.

If distance to $\mathcal E$ tends to zero, nearest equilibrium points exist and uniform continuity of $R$ implies residual convergence. Conversely, if residuals vanish but a subsequence stays a positive distance away, compactness supplies a convergent further subsequence. Its limit has residual zero by continuity, contradicting that distance.
\end{proof}
If an approximate value has uniform error $\epsilon$, each action gap has error at most $(1+\alpha)\epsilon$. Replacing approximate gaps by their positive parts changes (4) by at most $|X|(1+\alpha)\epsilon$, since the action probabilities sum to one at each state. This controls evaluation error in a diagnostic. It does not provide a universal linear bound of equilibrium distance by residual: continuity supplies a modulus only after the particular game is known.

Likewise, small policy changes between iterations are not equivalent to a small residual. A learner can stop updating at a nonoptimal policy because exploration ceases. Equation (4) tests stationarity and optimality directly, provided frozen optimal values are estimated accurately enough.
\section{Remaining problem and source boundary}
The exact cycle establishes that solving each frozen control problem is insufficient for convergence. The sampled theorem supplies a transparent special case with separate control of sample error and population movement. The residual result gives a correct target for numerical evaluation.

General noncontractive population games may cycle, possess multiple equilibria, or retain approximation and replay-buffer bias. The source's FP-DQN implementation needs its own stability argument and its actual approximation, exploration, averaging, and sample schedules. Those are not supplied by the one-state proof. The full catalogue problem therefore remains unresolved.

\section{Completion Estimate}
40\%

This is a post hoc, heuristic estimate for the full catalogue target.
The attempt proves constant-step cycling, almost-sure convergence of a specified sampled one-state benchmark, and a residual characterization of stationary equilibrium-set convergence. The intended noncontractive dynamic population models and actual FP-DQN approximation, exploration, replay, and averaging schedules still lack a convergence or failure theorem.

\begin{thebibliography}{9}
\bibitem{cbs} M. Cederle, S. Bolognani, and G. A. Susto (2026).
\emph{Towards Model-Free Learning in Dynamic Population Games: An Application to Karma Economies}, arXiv:2605.11042v1, sections 5 and 7.
The primary text separates numerical population convergence from its single-agent guarantees.
\url{https://arxiv.org/html/2605.11042v1}.
\end{thebibliography}

\end{document}
